\pdfoutput=1
\documentclass[copyright,creativecommons]{eptcs}
\providecommand{\event}{ICLP 2026}

% Math
\usepackage{amsmath}
\usepackage{amssymb,amsfonts}
\usepackage{mathtools}
\usepackage{bm}
\usepackage{relsize}

% Color (xcolor for \textcolor used by \udi, \claude)
\usepackage{xcolor}

% Graphics
\usepackage{graphicx}

% Misc
\usepackage{xspace}
\usepackage{cleveref}

% Algorithm packages
\usepackage{algorithm}
\usepackage[noend]{algpseudocode}

% Layout packages
\usepackage{multicol}
\usepackage{multirow}
\usepackage{verbatimbox}
\usepackage{wrapfig}
\usepackage{fancyvrb}

% Lists
\usepackage{enumitem}

% Theorem-like environments
\usepackage{amsthm}
\theoremstyle{definition}
\newtheorem{definition}{Definition}[section]
\newtheorem{example}[definition]{Example}
\newtheorem{observation}[definition]{Observation}
\theoremstyle{plain}
\newtheorem{theorem}[definition]{Theorem}
\newtheorem{proposition}[definition]{Proposition}
\newtheorem{lemma}[definition]{Lemma}
\newtheorem{corollary}[definition]{Corollary}
\theoremstyle{remark}
\newtheorem{remark}[definition]{Remark}

% Restatable theorems (must be after amsthm)
\usepackage{thmtools,thm-restate}

% Keywords environment (TLP supplied one; EPTCS does not)
\newenvironment{keywords}{\par\smallskip\noindent\textbf{Keywords:} }{}

% Set list spacing
\setlist{nosep, leftmargin=*}
\setlist{itemsep=1pt, topsep=3pt, leftmargin=*}

% Custom formatting commands
\newcommand{\mypara}[1]{\smallskip\noindent\textbf{#1.}}
\newcommand{\temph}[1]{\emph{#1}}

\newcommand{\remove}[1]{}
\newcommand{\udi}[1]{\textcolor{blue}{[Udi says: #1]}}
\newcommand{\claude}[1]{\textcolor{red}{[Claude: #1]}}

% Math commands
\newcommand{\calV}{\mathcal{V}}
\newcommand{\calG}{\mathcal{G}}
\newcommand{\calT}{\mathcal{T}}
\newcommand{\calF}{\mathcal{F}}
\newcommand{\calA}{\mathcal{A}}
\newcommand{\V}{\mathcal{V}}
\newcommand{\calS}{\mathcal{S}}
\newcommand{\calL}{\mathcal{L}}
\newcommand{\calP}{\mathcal{P}}
\newcommand{\calM}{\mathcal{M}}
\newcommand{\calE}{\mathcal{E}}
\newcommand{\calB}{\mathcal{B}}
\newcommand{\calC}{\mathcal{C}}
\newcommand{\calX}{\mathcal{X}}
\newcommand{\calD}{\mathcal{D}}

% Roman numeral abbreviations
\newcommand{\ia}{\textit{i}}
\newcommand{\ib}{\textit{ii}}
\newcommand{\ic}{\textit{iii}}
\newcommand{\id}{\textit{iv}}
\newcommand{\iie}{\textit{v}}
\newcommand{\iif}{\textit{vi}}
\newcommand{\iiv}{\textit{iv}}
\newcommand{\iv}{\textit{v}}

% Program counter for examples
\newcounter{pc}
\newcommand\spc{\addtocounter{pc}{1}\thepc}
\newcommand{\Program}[1]{\medskip\noindent\textbf{Program \spc: #1}\vspace{-5pt}}

% EPTCS volume cross-reference placeholder (per EPTCS instructions)
\providecommand{\thisvolume}[1]{this volume of EPTCS, Open Publishing Association}

\title{GLP: A Grassroots, Multiagent, Concurrent,\\ Logic Programming Language}
\author{Ehud Shapiro
\institute{London School of Economics\\ and Weizmann Institute of Science}
}

\hypersetup{
  bookmarksnumbered,
  pdftitle    = {GLP: A Grassroots, Multiagent, Concurrent, Logic Programming Language},
  pdfauthor   = {Ehud Shapiro},
  pdfsubject  = {EPTCS},
  pdfkeywords = {concurrent logic programming, grassroots platforms, operational semantics, multiagent transition systems}
}

\begin{document}

\maketitle

\begin{abstract}
Grassroots platforms are distributed systems with multiple instances that can (1) form and operate independently of each other and of any global resource other than the network, and (2) coalesce into ever larger instances, possibly resulting in a single global instance.

Here, we present Grassroots Logic Programs (GLP), a multiagent concurrent logic programming language designed for the implementation of grassroots platforms.  
We introduce the language  incrementally: We recall the standard operational semantics of logic programs; introduce the operational semantics of Concurrent (single-agent) GLP as a restriction of that of LP; recall the notion of multiagent transition systems and atomic transactions; introduce the operational semantics of multiagent GLP via a multiagent transition system specified via atomic transactions; and prove multiagent GLP to be grassroots. The accompanying programming example is the grassroots social graph---the infrastructure grassroots platform on which all others are based. 

With the mathematical foundations presented here: a workstation-based implementation of Concurrent GLP was developed by AI, based on the operational semantics of Concurrent GLP;  a distributed peer-to-peer smartphone-based implementation of multiagent GLP is being developed by AI, based on the operational semantics of multiagent GLP; a moded type system for GLP was implemented by AI, to facilitate the specification of GLP programs by human and AI designers, for their programming by AI; all reported in detail in companion papers.
\end{abstract}

\begin{keywords}
concurrent logic programming, grassroots platforms, operational semantics, multiagent transition systems
\end{keywords}

% Main sections
\input{sections/introduction}
\input{sections/concurrent-glp}
\input{sections/maglp}
\input{sections/grassroots}
\input{sections/related-work}
\section{Conclusion}\label{section:conclusion}

While concurrent logic programming and its powerful programming techniques have been known for four decades~\cite{shapiro1989family}, adoption has been hampered by their inaccessibility to the average programmer. Concurrent logic programming requires a substantial shift from sequential procedural or functional thinking to dataflow programming via partial bindings, with concurrent operational semantics in which communication occurs through binding paired logical variables which may contain further logic variables.  Our experience to date is that AI is highly effective and productive at programming from mathematical specifications when the target is a typed high-level language such as GLP, where types constrain the space of legal programs and provide a checkable specification at the human-AI interface~\cite{shapiro2026types}. Moreover, our preliminary experience suggests that the abstract nature of concurrent logic programming renders it a more powerful and productive language than mainstream languages for collaborative human-AI abstract specifications-based program development.

Constitutional governance of grassroots communities and federations builds on Constitutional Consensus~\cite{keidar2025constitutional,lewis2025morpheus} and digital social contracts~\cite{cardelli2020digital,shapiro2022foundations}, both realisable in GLP.  The GLP runtime and example programs are open-source at \url{https://github.com/EShapiro2/GLP}.


\bibliographystyle{eptcs}
\bibliography{bib}
\appendix

\input{sections/appendix-lp}
\input{sections/appendix-term-matching}
\input{sections/appendix-grassroots-defs}
\input{sections/appendix-proofs}
\input{sections/appendix-guards}
\input{sections/appendix-additional-techniques}
\input{sections/appendix-social-graph-walkthrough}
%==============================================================================
\section{Complete Social Graph Program}
\label{appendix:social-graph-complete}
%==============================================================================

This appendix presents the complete GLP social graph program that demonstrates cold-call befriending, text messaging, and friend-mediated introduction. The program includes three agents (Alice, Bob, Charlie) driven by actor scripts that execute a comprehensive scenario. A UI mediator sits between each agent and its user, translating between the agent's non-ground terms and the ground terms used at the user interface. We present the shared agent code, the UI mediator, the actors, and three boot variants: a single-process play with a network switch (cGLP), a multiagent boot with actors (maGLP), and a multiagent interactive UI boot (maGLP).

%------------------------------------------------------------------------------
\subsection{The Scenario}
\label{app:scenario}
%------------------------------------------------------------------------------

The play executes the following sequence:
\begin{enumerate}
\item Alice cold-calls Bob (Bob accepts) --- Alice and Bob become friends
\item Alice sends message to Bob: ``Hi Bob, this is Alice''
\item Bob cold-calls Charlie (Charlie accepts) --- Bob and Charlie become friends
\item Charlie sends message to Bob: ``Hi Bob, this is Charlie''
\item Bob introduces Alice to Charlie (both accept) --- Alice and Charlie become direct friends
\item Alice sends message to Charlie: ``Hi Charlie, this is Alice''
\item Charlie responds to Alice: ``Hi Alice, this is Charlie''
\end{enumerate}

%------------------------------------------------------------------------------
\subsection{Type Definitions}
\label{app:types}
%------------------------------------------------------------------------------

\begin{verbatim}
Response ::= accept(Channel) ; no.
Decision ::= yes ; no.
PendingValue ::= response(Response?) ; channel(Channel?) ; error.
OutputEntry ::= output(String, Stream?).
OutputsList ::= [] ; [OutputEntry|OutputsList].
\end{verbatim}

%------------------------------------------------------------------------------
\subsection{Channel Operations}
\label{app:channel-ops}
%------------------------------------------------------------------------------

\begin{verbatim}
send(X, ch(In, [X?|Out?]), ch(In?, Out)).
receive(X?, ch([X|In], Out?), ch(In?, Out)).
new_channel(ch(Xs?, Ys), ch(Ys?, Xs)).
\end{verbatim}

%------------------------------------------------------------------------------
\subsection{Stream Utilities}
\label{app:stream-utils}
%------------------------------------------------------------------------------

\begin{verbatim}
procedure merge(Stream?, Stream?, Stream).
merge([X|Xs], Ys, [X?|Zs?]) :- merge(Ys?, Xs?, Zs).
merge(Xs, [Y|Ys], [Y?|Zs?]) :- merge(Xs?, Ys?, Zs).
merge([], Ys, Ys?).
merge(Xs, [], Xs?).
\end{verbatim}

%------------------------------------------------------------------------------
\subsection{Outputs List Operations}
\label{app:outputs-ops}
%------------------------------------------------------------------------------

\begin{verbatim}
procedure lookup_send(String?, _?, OutputsList?, OutputsList).
lookup_send(Key, Msg, Outs, Outs1?) :-
    ground(Key?) |
    lookup_send_step(Key?, Msg?, Outs?, Outs1).

procedure lookup_send_step(String?, _?, OutputsList?, OutputsList).
lookup_send_step(Key, Msg,
    [output(K, [Msg?|Out1?])|Rest],
    [output(K?, Out1)|Rest?]) :-
        Key? =?= K? | true.
lookup_send_step(Key, Msg,
    [output(K, Out?)|Rest],
    [output(K?, Out)|Rest1?]) :-
        otherwise |
        lookup_send_step(Key?, Msg?, Rest?, Rest1).
lookup_send_step(_, _, [], []).

procedure inject_msg(Response?, Constant?, Constant?, Stream?, Stream).
inject_msg(Resp, Target, Id, Ys,
    [msg(Target?, Id?, response(Resp?))|Ys?]) :-
        known(Resp?) | true.
inject_msg(Resp, Target, Id, [Y|Ys], [Y?|Ys1?]) :-
    inject_msg(Resp?, Target?, Id?, Ys?, Ys1).

procedure add_output(String?, Stream, OutputsList?, OutputsList).
add_output(Name, Out?, Outs, [output(Name?, Out)|Outs?]).

procedure close_outputs(OutputsList?).
close_outputs([output(_, [])|Outs]) :- close_outputs(Outs?).
close_outputs([]).
\end{verbatim}

%------------------------------------------------------------------------------
\subsection{Response Handling}
\label{app:response-handling}
%------------------------------------------------------------------------------

\begin{verbatim}
procedure bind_response(Decision?, Constant?, Response,
    OutputsList?, OutputsList, Stream?, Stream).
bind_response(yes, From, accept(RetCh?),
    Outs, Outs1?, In, In1?) :-
        new_channel(RetCh, LocalCh) |
        handle_response(accept(LocalCh?), From?,
            Outs?, Outs1, In?, In1).
bind_response(no, _, no, Outs, Outs1?, In, In?) :-
    lookup_send('_user',
        msg(agent, '_user', rejected),
        Outs?, Outs1).

procedure handle_response(Response?, Constant?,
    OutputsList?, OutputsList, Stream?, Stream).
handle_response(accept(ch(FIn, FOut?)), From,
    Outs, Outs2?, In, In1?) :-
        ground(From?) |
        add_output(From?, FOut, Outs?, Outs1),
        lookup_send('_user',
            msg(agent, '_user', connected(From?)),
            Outs1?, Outs2),
        merge(In?, FIn?, In1).
handle_response(no, From, Outs, Outs1?, In, In?) :-
    ground(From?) |
    lookup_send('_user',
        msg(agent, '_user', rejected(From?)),
        Outs?, Outs1).
\end{verbatim}

%------------------------------------------------------------------------------
\subsection{The Agent}
\label{app:agent}
%------------------------------------------------------------------------------

The agent procedure is the main event loop handling all protocols. It takes four arguments: agent identity, user input stream, network input stream, and an outputs list mapping names to output streams.

\begin{verbatim}
procedure agent(Constant?, Stream?, Stream?, OutputsList?).

%% --- User messages ---

%% User initiates cold call
agent(Id, [msg('_user', Id1, connect(Target))|UserIn],
    NetIn, Outs) :-
        Id? =?= Id1?, ground(Target?) |
        lookup_send('_net',
            msg(Target?, intro(Id?, Resp)),
            Outs?, Outs1),
        inject_msg(Resp?, Target?, Id?,
            UserIn?, UserIn1),
        agent(Id?, UserIn1?, NetIn?, Outs1?).

%% User decision on cold-call
agent(Id, [msg('_user', Id1,
    decision(Dec, From, response(Resp?)))|UserIn],
    NetIn, Outs) :-
        Id? =?= Id1? |
        bind_response(Dec?, From?, Resp,
            Outs?, Outs1, NetIn?, NetIn1),
        agent(Id?, UserIn?, NetIn1?, Outs1?).

%% User sends text message
agent(Id, [msg('_user', Id1,
    send(Target, Text))|UserIn],
    NetIn, Outs) :-
        Id? =?= Id1?, ground(Target?) |
        lookup_send(Target?,
            msg(Id?, Target?, text(Text?)),
            Outs?, Outs1),
        agent(Id?, UserIn?, NetIn?, Outs1?).

%% User introduces P to Q
agent(Id, [msg('_user', Id1,
    introduce(P, Q))|UserIn],
    NetIn, Outs) :-
        Id? =?= Id1?, ground(P?), ground(Q?),
        new_channel(PQCh, QPCh) |
        lookup_send(P?,
            msg(Id?, P?, intro(Q?, QPCh?)),
            Outs?, Outs1),
        lookup_send(Q?,
            msg(Id?, Q?, intro(P?, PQCh?)),
            Outs1?, Outs2),
        agent(Id?, UserIn?, NetIn?, Outs2?).

%% User accepts friend introduction
agent(Id, [msg('_user', Id1,
    accept_intro(Other,
        channel(ch(FIn, FOut?))))|UserIn],
    NetIn, Outs) :-
        Id? =?= Id1?, ground(Other?) |
        add_output(Other?, FOut, Outs?, Outs1),
        lookup_send('_user',
            msg(agent, '_user', connected(Other?)),
            Outs1?, Outs2),
        merge(NetIn?, FIn?, NetIn1),
        agent(Id?, UserIn?, NetIn1?, Outs2?).

%% User rejects friend introduction
agent(Id, [msg('_user', Id1,
    reject_intro(_))|UserIn],
    NetIn, Outs) :-
        Id? =?= Id1? |
        agent(Id?, UserIn?, NetIn?, Outs?).

%% Response to cold-call (injected by inject_msg)
agent(Id, [msg(From, Id1, response(Resp))|UserIn],
    NetIn, Outs) :-
        Id? =?= Id1? |
        handle_response(Resp?, From?,
            Outs?, Outs1, NetIn?, NetIn1),
        agent(Id?, UserIn?, NetIn1?, Outs1?).

%% User catch-all
agent(Id, [_|UserIn], NetIn, Outs) :-
    ground(Id?), otherwise |
    agent(Id?, UserIn?, NetIn?, Outs?).

%% --- Network messages ---

%% Cold-call from network
agent(Id, UserIn,
    [msg(Id1, intro(From, Resp))|NetIn], Outs) :-
        Id? =?= Id1? |
        lookup_send('_user',
            msg(agent, '_user',
                befriend(From?, Resp?)),
            Outs?, Outs1),
        agent(Id?, UserIn?, NetIn?, Outs1?).

%% Response to cold-call from network
agent(Id, UserIn,
    [msg(From, Id1, response(Resp))|NetIn],
    Outs) :-
        Id? =?= Id1? |
        handle_response(Resp?, From?,
            Outs?, Outs1, NetIn?, NetIn1),
        agent(Id?, UserIn?, NetIn1?, Outs1?).

%% Text from friend
agent(Id, UserIn,
    [msg(From, Id1, text(Text))|NetIn], Outs) :-
        Id? =?= Id1? |
        lookup_send('_user',
            msg(agent, '_user',
                received(From?, Text?)),
            Outs?, Outs1),
        agent(Id?, UserIn?, NetIn?, Outs1?).

%% Introduction from friend
agent(Id, UserIn,
    [msg(From, Id1, intro(Other, Ch))|NetIn],
    Outs) :-
        Id? =?= Id1?, ground(Other?) |
        lookup_send('_user',
            msg(agent, '_user',
                befriend_intro(From?, Other?, Ch?)),
            Outs?, Outs1),
        agent(Id?, UserIn?, NetIn?, Outs1?).

%% Network catch-all
agent(Id, UserIn, [_|NetIn], Outs) :-
    ground(Id?), otherwise |
    agent(Id?, UserIn?, NetIn?, Outs?).

%% --- Termination ---

agent(_, [], _, Outs) :- close_outputs(Outs?).
agent(_, _, [], Outs) :- close_outputs(Outs?).
\end{verbatim}

%------------------------------------------------------------------------------
\subsection{Network Switch}
\label{app:network-switch}
%------------------------------------------------------------------------------

The three-way network switch routes messages between Alice, Bob, and Charlie. It handles both 2-argument cold-call messages and 3-argument friend-to-friend messages.

\begin{verbatim}
procedure network3(Channel?, Channel?, Channel?).

%% --- Cold-call routing (2-arg msg format) ---

%% Alice cold-calls Bob
network3(ch([msg(bob, X)|AliceIn], AliceOut?),
         ch(BobIn, [msg(bob, X?)|BobOut?]),
         ch(CharlieIn, CharlieOut?)) :-
    network3(ch(AliceIn?, AliceOut),
        ch(BobIn?, BobOut),
        ch(CharlieIn?, CharlieOut)).

%% Alice cold-calls Charlie
network3(
    ch([msg(charlie, X)|AliceIn], AliceOut?),
    ch(BobIn, BobOut?),
    ch(CharlieIn,
        [msg(charlie, X?)|CharlieOut?])) :-
    network3(ch(AliceIn?, AliceOut),
        ch(BobIn?, BobOut),
        ch(CharlieIn?, CharlieOut)).

%% Bob cold-calls Alice
network3(
    ch(AliceIn,
        [msg(alice, X?)|AliceOut?]),
    ch([msg(alice, X)|BobIn], BobOut?),
    ch(CharlieIn, CharlieOut?)) :-
    network3(ch(AliceIn?, AliceOut),
        ch(BobIn?, BobOut),
        ch(CharlieIn?, CharlieOut)).

%% Bob cold-calls Charlie
network3(ch(AliceIn, AliceOut?),
    ch([msg(charlie, X)|BobIn], BobOut?),
    ch(CharlieIn,
        [msg(charlie, X?)|CharlieOut?])) :-
    network3(ch(AliceIn?, AliceOut),
        ch(BobIn?, BobOut),
        ch(CharlieIn?, CharlieOut)).

%% Charlie cold-calls Alice
network3(
    ch(AliceIn,
        [msg(alice, X?)|AliceOut?]),
    ch(BobIn, BobOut?),
    ch([msg(alice, X)|CharlieIn],
        CharlieOut?)) :-
    network3(ch(AliceIn?, AliceOut),
        ch(BobIn?, BobOut),
        ch(CharlieIn?, CharlieOut)).

%% Charlie cold-calls Bob
network3(ch(AliceIn, AliceOut?),
    ch(BobIn, [msg(bob, X?)|BobOut?]),
    ch([msg(bob, X)|CharlieIn],
        CharlieOut?)) :-
    network3(ch(AliceIn?, AliceOut),
        ch(BobIn?, BobOut),
        ch(CharlieIn?, CharlieOut)).

%% --- Friend-to-friend routing (3-arg msg) ---

%% Alice -> Bob
network3(
    ch([msg(alice, bob, X)|AliceIn],
        AliceOut?),
    ch(BobIn,
        [msg(alice, bob, X?)|BobOut?]),
    ch(CharlieIn, CharlieOut?)) :-
    network3(ch(AliceIn?, AliceOut),
        ch(BobIn?, BobOut),
        ch(CharlieIn?, CharlieOut)).

%% Alice -> Charlie
network3(
    ch([msg(alice, charlie, X)|AliceIn],
        AliceOut?),
    ch(BobIn, BobOut?),
    ch(CharlieIn,
        [msg(alice, charlie, X?)
         |CharlieOut?])) :-
    network3(ch(AliceIn?, AliceOut),
        ch(BobIn?, BobOut),
        ch(CharlieIn?, CharlieOut)).

%% Bob -> Alice
network3(
    ch(AliceIn,
        [msg(bob, alice, X?)|AliceOut?]),
    ch([msg(bob, alice, X)|BobIn],
        BobOut?),
    ch(CharlieIn, CharlieOut?)) :-
    network3(ch(AliceIn?, AliceOut),
        ch(BobIn?, BobOut),
        ch(CharlieIn?, CharlieOut)).

%% Bob -> Charlie
network3(ch(AliceIn, AliceOut?),
    ch([msg(bob, charlie, X)|BobIn],
        BobOut?),
    ch(CharlieIn,
        [msg(bob, charlie, X?)
         |CharlieOut?])) :-
    network3(ch(AliceIn?, AliceOut),
        ch(BobIn?, BobOut),
        ch(CharlieIn?, CharlieOut)).

%% Charlie -> Alice
network3(
    ch(AliceIn,
        [msg(charlie, alice, X?)
         |AliceOut?]),
    ch(BobIn, BobOut?),
    ch([msg(charlie, alice, X)|CharlieIn],
        CharlieOut?)) :-
    network3(ch(AliceIn?, AliceOut),
        ch(BobIn?, BobOut),
        ch(CharlieIn?, CharlieOut)).

%% Charlie -> Bob
network3(ch(AliceIn, AliceOut?),
    ch(BobIn,
        [msg(charlie, bob, X?)|BobOut?]),
    ch([msg(charlie, bob, X)|CharlieIn],
        CharlieOut?)) :-
    network3(ch(AliceIn?, AliceOut),
        ch(BobIn?, BobOut),
        ch(CharlieIn?, CharlieOut)).

%% Termination
network3(ch([], []), ch([], []), ch([], [])).
\end{verbatim}

%------------------------------------------------------------------------------
\subsection{UI Mediator}
\label{app:ui-mediator}
%------------------------------------------------------------------------------

The UI mediator sits between \verb|agent/4| and the user (whether a human or an actor). Agent output may contain unbound variables (e.g., \verb|befriend(alice, Resp)|); the mediator replaces these with ground request identifiers \verb|req(N)|, stores the original variables in a pending list, and retrieves them when the user responds. All communication between the mediator and the user is in ground terms.

\begin{verbatim}
ReqId       ::= req(Constant).
PendingValue ::= response(Response?)
               ; channel(Channel?) ; error.
PendingEntry ::= pending(ReqId, PendingValue).
PendingList  ::= [] ; [PendingEntry|PendingList].

procedure ui_mediator(Constant?, Channel?,
    Channel?, PendingList?, Constant?).

%% --- Agent-to-user: non-ground ---

%% Cold-call befriend: store response variable
ui_mediator(Id, AgentCh, UserCh, Ps, N) :-
    receive(msg(agent, '_user',
                befriend(From, Resp?)),
            AgentCh?, AgentCh1),
    ground(From?) |
    send(befriend(From?, req(N?)),
         UserCh?, UserCh1),
    N1 := N? + 1,
    ui_mediator(Id?, AgentCh1?, UserCh1?,
        [pending(req(N?), response(Resp))
         | Ps?], N1?).

%% Introduction: store channel variable
ui_mediator(Id, AgentCh, UserCh, Ps, N) :-
    receive(msg(agent, '_user',
                befriend_intro(From, Other, Ch?)),
            AgentCh?, AgentCh1),
    ground(From?), ground(Other?) |
    send(befriend_intro(From?, Other?, req(N?)),
         UserCh?, UserCh1),
    N1 := N? + 1,
    ui_mediator(Id?, AgentCh1?, UserCh1?,
        [pending(req(N?), channel(Ch))
         | Ps?], N1?).

%% --- Agent-to-user: ground (pass through) ---

ui_mediator(Id, AgentCh, UserCh, Ps, N) :-
    receive(msg(agent, '_user', connected(Who)),
            AgentCh?, AgentCh1),
    ground(Who?) |
    send(connected(Who?), UserCh?, UserCh1),
    ui_mediator(Id?, AgentCh1?, UserCh1?,
        Ps?, N?).

ui_mediator(Id, AgentCh, UserCh, Ps, N) :-
    receive(msg(agent, '_user', rejected),
            AgentCh?, AgentCh1) |
    send(rejected, UserCh?, UserCh1),
    ui_mediator(Id?, AgentCh1?, UserCh1?,
        Ps?, N?).

ui_mediator(Id, AgentCh, UserCh, Ps, N) :-
    receive(msg(agent, '_user', rejected(Who)),
            AgentCh?, AgentCh1),
    ground(Who?) |
    send(rejected(Who?), UserCh?, UserCh1),
    ui_mediator(Id?, AgentCh1?, UserCh1?,
        Ps?, N?).

ui_mediator(Id, AgentCh, UserCh, Ps, N) :-
    receive(msg(agent, '_user',
                received(From, Text)),
            AgentCh?, AgentCh1),
    ground(From?), ground(Text?) |
    send(received(From?, Text?),
         UserCh?, UserCh1),
    ui_mediator(Id?, AgentCh1?, UserCh1?,
        Ps?, N?).

%% --- User-to-agent: decision ---

ui_mediator(Id, AgentCh, UserCh, Ps, N) :-
    receive(decision(Dec, From, ReqId),
            UserCh?, UserCh1),
    ground(Id?), ground(Dec?), ground(From?),
    ground(ReqId?) |
    lookup_pending(ReqId?, Resp, Ps?, Ps1),
    send(msg('_user', Id?,
             decision(Dec?, From?, Resp?)),
         AgentCh?, AgentCh1),
    ui_mediator(Id?, AgentCh1?, UserCh1?,
        Ps1?, N?).

%% --- User-to-agent: accept/reject intro ---

ui_mediator(Id, AgentCh, UserCh, Ps, N) :-
    receive(accept_intro(Other, ReqId),
            UserCh?, UserCh1),
    ground(Id?), ground(Other?),
    ground(ReqId?) |
    lookup_pending(ReqId?, Ch, Ps?, Ps1),
    send(msg('_user', Id?,
             accept_intro(Other?, Ch?)),
         AgentCh?, AgentCh1),
    ui_mediator(Id?, AgentCh1?, UserCh1?,
        Ps1?, N?).

ui_mediator(Id, AgentCh, UserCh, Ps, N) :-
    receive(reject_intro(Other),
            UserCh?, UserCh1),
    ground(Id?), ground(Other?) |
    send(msg('_user', Id?,
             reject_intro(Other?)),
         AgentCh?, AgentCh1),
    ui_mediator(Id?, AgentCh1?, UserCh1?,
        Ps?, N?).

%% --- User-to-agent: pass-through commands ---

ui_mediator(Id, AgentCh, UserCh, Ps, N) :-
    receive(connect(Target),
            UserCh?, UserCh1),
    ground(Id?), ground(Target?) |
    send(msg('_user', Id?,
             connect(Target?)),
         AgentCh?, AgentCh1),
    ui_mediator(Id?, AgentCh1?, UserCh1?,
        Ps?, N?).

ui_mediator(Id, AgentCh, UserCh, Ps, N) :-
    receive(send(Target, Text),
            UserCh?, UserCh1),
    ground(Id?), ground(Target?),
    ground(Text?) |
    send(msg('_user', Id?,
             send(Target?, Text?)),
         AgentCh?, AgentCh1),
    ui_mediator(Id?, AgentCh1?, UserCh1?,
        Ps?, N?).

ui_mediator(Id, AgentCh, UserCh, Ps, N) :-
    receive(introduce(P, Q),
            UserCh?, UserCh1),
    ground(Id?), ground(P?), ground(Q?) |
    send(msg('_user', Id?,
             introduce(P?, Q?)),
         AgentCh?, AgentCh1),
    ui_mediator(Id?, AgentCh1?, UserCh1?,
        Ps?, N?).

%% Termination
ui_mediator(_, ch([], []),
    ch([], []), _, _).

%% --- Pending list management ---

procedure lookup_pending(ReqId?,
    PendingValue, PendingList?, PendingList).
lookup_pending(ReqId, Val?,
    [pending(Id, Val) | Rest], Rest?) :-
    ReqId? =?= Id? | true.
lookup_pending(ReqId, Val?,
    [P | Rest], [P? | Rest1?]) :-
    otherwise |
    lookup_pending(ReqId?, Val, Rest?, Rest1).
lookup_pending(_, error, [], []).
\end{verbatim}

%------------------------------------------------------------------------------
\subsection{UI Actors}
\label{app:actors}
%------------------------------------------------------------------------------

Each actor is a GLP procedure that implements a state machine, communicating with the UI mediator in ground terms---the same protocol a human uses in the interactive UI. The actor's state is encoded in the procedure name (e.g., \verb|alice_ui_wait_bob_connected|, \verb|bob_ui_wait_charlie_msg|), and transitions occur via recursive calls to different procedures.

\subsubsection{Alice's Actor}

Alice's script: (1) send \verb|connect(bob)|, (2) wait for \verb|connected(bob)|, (3) send greeting to Bob, (4) wait for \verb|befriend_intro| and accept it, (5) wait for \verb|connected(charlie)|, (6) send greeting to Charlie, (7) wait for Charlie's reply.

\begin{verbatim}
procedure alice_ui_actor(Channel?).
alice_ui_actor(ch(In,
    [connect(bob)|Out?])) :-
        alice_ui_wait_bob_connected(In?, Out).

procedure alice_ui_wait_bob_connected(
    Stream?, Stream).
alice_ui_wait_bob_connected(
    [connected(bob)|In],
    [send(bob, 'Hi Bob, this is Alice')
     |Out?]) :-
        alice_ui_wait_intro(In?, Out).
alice_ui_wait_bob_connected([_|In], Out?) :-
    otherwise |
    alice_ui_wait_bob_connected(In?, Out).
alice_ui_wait_bob_connected([], []).

procedure alice_ui_wait_intro(
    Stream?, Stream).
alice_ui_wait_intro(
    [befriend_intro(bob, charlie, ReqId)|In],
    [accept_intro(charlie, ReqId?)|Out?]) :-
        ground(ReqId?) |
        alice_ui_wait_charlie_connected(
            In?, Out).
alice_ui_wait_intro([_|In], Out?) :-
    otherwise |
    alice_ui_wait_intro(In?, Out).
alice_ui_wait_intro([], []).

procedure alice_ui_wait_charlie_connected(
    Stream?, Stream).
alice_ui_wait_charlie_connected(
    [connected(charlie)|In],
    [send(charlie,
        'Hi Charlie, this is Alice')
     |Out?]) :-
        alice_ui_wait_charlie_reply(In?, Out).
alice_ui_wait_charlie_connected(
    [_|In], Out?) :-
        otherwise |
        alice_ui_wait_charlie_connected(
            In?, Out).
alice_ui_wait_charlie_connected([], []).

procedure alice_ui_wait_charlie_reply(
    Stream?, Stream).
alice_ui_wait_charlie_reply(
    [received(charlie, _)|_], []).
alice_ui_wait_charlie_reply([_|In], Out?) :-
    otherwise |
    alice_ui_wait_charlie_reply(In?, Out).
alice_ui_wait_charlie_reply([], []).
\end{verbatim}

\subsubsection{Bob's Actor}

Bob's script: (1) wait for \verb|befriend(alice, req(N))| and accept, (2) wait for Alice's message, (3) send \verb|connect(charlie)|, (4) wait for connection, (5) wait for Charlie's message, (6) introduce Alice to Charlie.

\begin{verbatim}
procedure bob_ui_actor(Channel?).
bob_ui_actor(ch(
    [befriend(alice, ReqId)|In],
    [decision(yes, alice, ReqId?)|Out?])) :-
        ground(ReqId?) |
        bob_ui_wait_alice_msg(In?, Out).
bob_ui_actor(ch([_|In], Out?)) :-
    otherwise | bob_ui_actor(ch(In?, Out)).

procedure bob_ui_wait_alice_msg(
    Stream?, Stream).
bob_ui_wait_alice_msg(
    [connected(_)|In], Out?) :-
        bob_ui_wait_alice_msg(In?, Out).
bob_ui_wait_alice_msg(
    [received(alice, _)|In],
    [connect(charlie)|Out?]) :-
        bob_ui_wait_charlie_connected(
            In?, Out).
bob_ui_wait_alice_msg([_|In], Out?) :-
    otherwise |
    bob_ui_wait_alice_msg(In?, Out).
bob_ui_wait_alice_msg([], []).

procedure bob_ui_wait_charlie_connected(
    Stream?, Stream).
bob_ui_wait_charlie_connected(
    [connected(charlie)|In], Out?) :-
        bob_ui_wait_charlie_msg(In?, Out).
bob_ui_wait_charlie_connected(
    [_|In], Out?) :-
        otherwise |
        bob_ui_wait_charlie_connected(
            In?, Out).
bob_ui_wait_charlie_connected([], []).

procedure bob_ui_wait_charlie_msg(
    Stream?, Stream).
bob_ui_wait_charlie_msg(
    [received(charlie, _)|_],
    [introduce(alice, charlie)]).
bob_ui_wait_charlie_msg([_|In], Out?) :-
    otherwise |
    bob_ui_wait_charlie_msg(In?, Out).
bob_ui_wait_charlie_msg([], []).
\end{verbatim}

\subsubsection{Charlie's Actor}

Charlie's script: (1) wait for \verb|befriend(bob, req(N))| and accept, (2) send greeting to Bob, (3) wait for \verb|befriend_intro| and accept, (4) wait for Alice's message, (5) send reply to Alice.

\begin{verbatim}
procedure charlie_ui_actor(Channel?).
charlie_ui_actor(ch(
    [befriend(bob, ReqId)|In],
    [decision(yes, bob, ReqId?),
     send(bob, 'Hi Bob, this is Charlie')
     |Out?])) :-
        ground(ReqId?) |
        charlie_ui_wait_intro(In?, Out).
charlie_ui_actor(ch([_|In], Out?)) :-
    otherwise |
    charlie_ui_actor(ch(In?, Out)).

procedure charlie_ui_wait_intro(
    Stream?, Stream).
charlie_ui_wait_intro(
    [befriend_intro(bob, alice, ReqId)|In],
    [accept_intro(alice, ReqId?)|Out?]) :-
        ground(ReqId?) |
        charlie_ui_wait_alice_msg(In?, Out).
charlie_ui_wait_intro([_|In], Out?) :-
    otherwise |
    charlie_ui_wait_intro(In?, Out).
charlie_ui_wait_intro([], []).

procedure charlie_ui_wait_alice_msg(
    Stream?, Stream).
charlie_ui_wait_alice_msg(
    [received(alice, _)|_],
    [send(alice,
        'Hi Alice, this is Charlie')]).
charlie_ui_wait_alice_msg([_|In], Out?) :-
    otherwise |
    charlie_ui_wait_alice_msg(In?, Out).
charlie_ui_wait_alice_msg([], []).
\end{verbatim}

%------------------------------------------------------------------------------
\subsection{The Play (cGLP)}
\label{app:play}
%------------------------------------------------------------------------------

The play initializes three agents, each connected to the network switch via the UI mediator. All agents run in a single process.

\begin{verbatim}
procedure play.
play :-
    network3(ch(AliceNetOut?, AliceNetIn),
             ch(BobNetOut?, BobNetIn),
             ch(CharlieNetOut?,
                CharlieNetIn)),

    %% Alice: agent <-> mediator <-> actor
    agent(alice, AliceAgentIn?,
        AliceNetIn?,
        [output('_user', AliceAgentToUser),
         output('_net', AliceNetOut)]),
    ui_mediator(alice,
        ch(AliceAgentToUser?, AliceAgentIn),
        ch(AliceUserOut?, AliceUserIn),
        [], 1),
    alice_ui_actor(
        ch(AliceUserIn?, AliceUserOut)),

    %% Bob: agent <-> mediator <-> actor
    agent(bob, BobAgentIn?,
        BobNetIn?,
        [output('_user', BobAgentToUser),
         output('_net', BobNetOut)]),
    ui_mediator(bob,
        ch(BobAgentToUser?, BobAgentIn),
        ch(BobUserOut?, BobUserIn),
        [], 1),
    bob_ui_actor(
        ch(BobUserIn?, BobUserOut)),

    %% Charlie: agent <-> mediator <-> actor
    agent(charlie, CharlieAgentIn?,
        CharlieNetIn?,
        [output('_user', CharlieAgentToUser),
         output('_net', CharlieNetOut)]),
    ui_mediator(charlie,
        ch(CharlieAgentToUser?,
           CharlieAgentIn),
        ch(CharlieUserOut?,
           CharlieUserIn),
        [], 1),
    charlie_ui_actor(
        ch(CharlieUserIn?,
           CharlieUserOut)).
\end{verbatim}

The play terminates successfully when all streams are closed and all agents reach their terminal states. The network switch simulates the Cold-call transaction introduced in Section~\ref{sec:maglp}.

%------------------------------------------------------------------------------
\subsection{Multiagent Boot (maGLP with Actors)}
\label{app:maglp-boot}
%------------------------------------------------------------------------------

In the multiagent setting, each agent runs on a separate isolate. The \verb|boot| procedure spawns agents using the \verb|@agent| syntax, which places each goal on the named agent's isolate. The \verb|send_to_net| system predicate returns the agent's network output stream. No network switch is needed---maGLP's Cold-call transaction handles inter-agent communication directly.

\begin{verbatim}
procedure boot.
boot :-
    agent_init(alice, _)@alice,
    agent_init(bob, _)@bob,
    agent_init(charlie, _)@charlie.

procedure agent_init(Constant?, Stream?).
agent_init(Id, NetIn) :-
    ground(Id?) |
    send_to_net(NetOut?),
    agent(Id?, AgentIn?, NetIn?,
        [output('_user', AgentToUser),
         output('_net', NetOut)]),
    ui_mediator(Id?,
        ch(AgentToUser?, AgentIn),
        ch(UserOut?, UserIn), [], 1),
    ui_actor(Id?,
        ch(UserIn?, UserOut)).

procedure ui_actor(_?, Channel?).
ui_actor(alice, Ch) :-
    alice_ui_actor(Ch?).
ui_actor(bob, Ch) :-
    bob_ui_actor(Ch?).
ui_actor(charlie, Ch) :-
    charlie_ui_actor(Ch?).
\end{verbatim}

%------------------------------------------------------------------------------
\subsection{Multiagent UI Boot (maGLP Interactive)}
\label{app:maglp-ui-boot}
%------------------------------------------------------------------------------

In the interactive UI variant, each agent runs in its own window. The user interacts by typing ground GLP terms (e.g., \verb|connect(bob)|, \verb|decision(yes, alice, req(1))|, \verb|introduce(alice, charlie)|). The \verb|send_to_user| system predicate outputs ground terms from the mediator to the UI.

\begin{verbatim}
agent_init(Id, UserIn, NetIn) :-
    ground(Id?) |
    send_to_net(NetOut?),
    agent(Id?, AgentIn?, NetIn?,
        [output('_user', AgentToUser),
         output('_net', NetOut)]),
    ui_mediator(Id?,
        ch(AgentToUser?, AgentIn),
        ch(UserIn?, UserOut), [], 1),
    send_to_user(UserOut?).
\end{verbatim}


\end{document}
